% ============================================================
% DocuForge — Manuscript to PPT: Beamer Preamble
% 此文件由流水线自动拼接在 LLM 生成的 frame 内容之前。
% LLM 只需输出 \begin{frame} ... \end{frame} 和 \end{document}。
% ============================================================

\documentclass[10pt, aspectratio=169, t]{beamer}

% ---------- 主题 ----------
\usetheme[progressbar=frametitle, numbering=fraction, sectionpage=none]{metropolis}

% 自定义配色 —— 深蓝+灰白色调，学术感
\definecolor{PrimaryBlue}{HTML}{2C3E6B}
\definecolor{AccentTeal}{HTML}{1A8A8A}
\definecolor{LightBg}{HTML}{F5F5F5}
\definecolor{DarkText}{HTML}{2D2D2D}
\definecolor{MutedGray}{HTML}{6B7280}
\definecolor{HighlightRed}{HTML}{C0392B}

\setbeamercolor{palette primary}{bg=PrimaryBlue, fg=white}
\setbeamercolor{frametitle}{bg=PrimaryBlue, fg=white}
\setbeamercolor{progress bar}{fg=AccentTeal}
\setbeamercolor{title separator}{fg=AccentTeal}
\setbeamercolor{alerted text}{fg=HighlightRed}
\setbeamercolor{block title}{bg=PrimaryBlue!15, fg=PrimaryBlue}
\setbeamercolor{block body}{bg=LightBg}
\setbeamercolor{normal text}{fg=DarkText}

% ---------- 字体 ----------
\usepackage{fontspec}
\setsansfont{Noto Sans}
\usepackage{xeCJK}
\setCJKsansfont{Noto Sans CJK SC}
\setCJKmainfont{Noto Sans CJK SC}

% ---------- 数学 ----------
\usepackage{amsmath, amssymb, mathtools, bm}

% ---------- 绘图与代码 ----------
\usepackage{tikz}
\usetikzlibrary{arrows.meta, positioning, calc, shapes.geometric, decorations.pathreplacing}
\usepackage{algorithm2e}
\SetAlgorithmName{算法}{算法}{算法列表}

% ---------- 其他 ----------
\usepackage{booktabs}
\usepackage{graphicx}
\usepackage{hyperref}

% 紧凑间距
\setlength{\parskip}{0.3em}
\setbeamertemplate{itemize items}[circle]
\setbeamertemplate{enumerate items}[default]

% 脚注样式
\setbeamerfont{footnote}{size=\tiny}

% ============================================================
% ↓↓↓ 以下由 LLM 生成的内容会被拼接在此处 ↓↓↓
% ============================================================
